%%%%%%%%%%  scan idx 47  |  folio 36 (continuación)  %%%%%%%%%%

\underline{Definición 1.4.} (1) \underline{Sea $R$ un objeto de
$\underline C_\Lambda$. Una pequeña extensión de $R$ en
$\underline C_\Lambda$ es un morfismo sobreyectivo
$R'\xrightarrow{f}R$ en $\underline C_\Lambda$ tal que
$(\ker f)^2=0$ y $\ker f\simeq k$. En otras palabras,
$R'\xrightarrow{f}R$ es una pequeña extensión de $R$ si}
\[
0\longrightarrow k\xrightarrow{t}R'\xrightarrow{f}R
 \longrightarrow0
\]
\underline{es exacta para algún $t\in R'$ con $t^2=0$.}

(2) \underline{Sea $\underline A$ un grupoide cofibrado sobre
$\underline C_\Lambda$, y sea $a$ un objeto de $\underline A$.
Una pequeña prolongación de $a$ es una flecha
$a'\xrightarrow{\alpha}a$ en $\underline A$ tal que $P(\alpha)$ es
una pequeña extensión en $\underline C_\Lambda$. Una flecha
$a'\xrightarrow{\alpha}a$ se llamará versal para las pequeñas
prolongaciones de $a$ si, para toda pequeña prolongación
$a'_1\longrightarrow a$, existe un diagrama conmutativo}
\[
\begin{tikzcd}[row sep=large,column sep=large]
a' \arrow[rr] \arrow[dr,"\alpha"'] & & a'_1 \arrow[dl] \\
& a &
\end{tikzcd}
\]

\underline{Proposición 1.5. Sea $\underline A$ un grupoide cofibrado
semihomogéneo sobre $\underline C_\Lambda$.}

(1) \underline{Sean}
\[
\begin{tikzcd}[column sep=large]
& R \\
R\times_k k[\varepsilon]
 \arrow[ur,"\pi_1"] \arrow[dr,"\pi_2"'] & \\
& k[\varepsilon]
\end{tikzcd}
\]
\underline{los morfismos de proyección en $\underline C_\Lambda$, y
sea $a'\xrightarrow{\alpha}a$ un} $\pi_1$-\underline{morfismo. Entonces existe un
morfismo $a\xrightarrow{\beta}a'$ tal que $\alpha\beta=1_a$ si y solo
si existe un morfismo $a\longrightarrow(\pi_2)_*a'$.}

%%%%%%%%%%  scan idx 48  |  folio 37  %%%%%%%%%%

(2) \underline{Sea $R'\xrightarrow{f}R$ una pequeña extensión en
$\underline C_\Lambda$, y consideremos un diagrama conmutativo}
\[
\begin{tikzcd}[row sep=large,column sep=small]
& a'' \arrow[dl,"\alpha'"'] \arrow[dr,"\alpha'_1"]
  & & & & & R'\times_RR' \arrow[dl] \arrow[dr] & \\
a' \arrow[dr,"\alpha"'] & & \sisi{a'}{a'_1} \arrow[dl,"\alpha_1"]
  & & \underline{\text{situado sobre}} & R' \arrow[dr] & & R' \arrow[dl] \\
& a & & & & & R &
\end{tikzcd}
\]
\underline{Definamos el morfismo canónico
$\widetilde f:R'\times_RR'\longrightarrow k[\ker f]$ mediante}
\[
 \widetilde f(r'_1,r'_2)=r'_1(0)+(r'_1-r'_2).
\]
\underline{Si existe un morfismo
$a'\longrightarrow(\widetilde f)_*(a'')$, entonces existe un morfismo
$\gamma:a'\longrightarrow \sisi{a'}{a'_1}$ tal que
$\alpha=\alpha_1\gamma$.}

\underline{Demostración.} (1) La parte «solo si» es clara. Supongamos,
recíprocamente, que existe un morfismo
$\varphi:a\longrightarrow(\pi_2)_*a'$. Pongamos
\[
 h:R\longrightarrow R\times_k k[\varepsilon],
 \qquad h=1\times P(\varphi).
\]
El compuesto $\pi_1h$ es la identidad. Obtenemos, por tanto, morfismos
\[
 a\mathrel{\mathop{\rightleftarrows}^{\gamma}_{\alpha_1}}h_*a,
 \qquad \alpha_1\gamma=1_a,
 \qquad P(\gamma)=h,
 \qquad P(\alpha_1)=\pi_1.
\]
Por otra parte, $\pi_2h=P(\varphi)$, de modo que existe un
$\pi_2$-morfismo
\[
 h_*a\longrightarrow(\pi_2)_*a'.
\]
Así, $(\pi_2)_*(h_*a)=(\pi_2)_*a'$, y 1.2(2) proporciona un
$\underline C_\Lambda$-morfismo $u:h_*a\longrightarrow a'$ tal que
$\alpha u=\alpha_1$. Definamos $\beta=u\gamma$. Entonces
\[
 \alpha\beta=\alpha u\gamma=\alpha_1\gamma=1_a.
\]

(2) El morfismo
\[
 1\times\widetilde f:R'\times_RR'
 \longrightarrow R'\times_k k[\ker f]
\]
es un isomorfismo compatible con las proyecciones sobre el primer
factor, y $\pi_2(1\times\widetilde f)=\widetilde f$, donde
$\pi_2:R'\times_k k[\ker f]\longrightarrow k[\ker f]$ es la segunda
proyección. Puesto que $R'\xrightarrow fR$ es una pequeña extensión,
$k[\ker f]\simeq k[\varepsilon]$. Por consiguiente, si existe un
morfismo
\[
 a'\longrightarrow(\widetilde f)_*a''
   =(\pi_2)_*\bigl((1\times\widetilde f)_*a''\bigr),
\]
la parte (1) proporciona un morfismo $\gamma':a'\longrightarrow a''$
tal que $\alpha'\gamma'=1_{a'}$. Definamos
$\gamma=\alpha'_1\gamma':a'\longrightarrow \sisi{a'}{a'_1}$. Entonces
$\alpha_1\gamma=\alpha$.

%%%%%%%%%%  scan idx 49  |  folio 38  %%%%%%%%%%

\underline{Corolario 1.6. Sea $\underline A$ un grupoide cofibrado
semihomogéneo sobre $\underline C_\Lambda$, y consideremos un diagrama}
\[
\begin{tikzcd}[column sep=large]
a' \arrow[dr,"\alpha"'] & & a_1 \arrow[dl,"\alpha_1"] \\
& a &
\end{tikzcd}
\]
\underline{en $\underline A$, donde $\alpha_1$ es una pequeña
prolongación. Supongamos que $a$ es versal para las pequeñas
prolongaciones de $e$. Entonces existe un morfismo
$\rho:a'\sisi{\relbar\joinrel\relbar\joinrel\relbar}{\longrightarrow} a_1$ tal que $\alpha=\alpha_1\rho$ si y solo
si existe un morfismo $f:P(a')\longrightarrow P(a_1)$ tal que
$P(\alpha)=P(\alpha_1)f$ en $\underline C_\Lambda$.}

\underline{Demostración.} Supongamos que
$P(\alpha)=P(\alpha_1)f$, donde
$f:P(a')\longrightarrow P(a_1)$. Sustituyendo $a'$ por su imagen
directa mediante $f$, podemos y vamos a suponer que
$P(\alpha)=P(\alpha_1)$. Pongamos $R'=P(a_1)$ y $R=P(a)$. Como
$\underline A$ es semihomogéneo, obtenemos un diagrama conmutativo
\[
\begin{tikzcd}[row sep=large,column sep=small]
& a'' \arrow[dl,"\beta"'] \arrow[dr,"\beta_1"]
  & & & & & R'\times_RR' \arrow[dl] \arrow[dr] & \\
a' \arrow[dr,"\alpha"'] & & a_1 \arrow[dl,"\alpha_1"]
  & & \text{situado sobre} & R' \arrow[dr] & & R' \arrow[dl] \\
& a & & & & & R &
\end{tikzcd}
\]
Como $a$ es versal para las pequeñas prolongaciones de $e$, también lo
es $a'$; la afirmación se deduce entonces de 1.5(2).
